% Standalone independent audit of exact P08 v1.0, 2026-10-03.
% MB FOUNDATION VERDICT: VERIFIED
% MB SPECTRAL VERDICT: VERIFIED
% SAME-TUPLE COMPARISON VERDICT: MINOR REVISION
% PARITY/OBSTRUCTION VERDICT: VERIFIED
% OVERALL VERDICT: MINOR REVISION
% Required correction E01: the universal witness selects MB representative 2,
% not 1, at n=8. P08 source and all lifecycle metadata remain unchanged.
\documentclass[11pt]{article}
\usepackage[T1]{fontenc}
\usepackage[utf8]{inputenc}
\usepackage{lmodern,amsmath,amssymb,amsthm}
\usepackage[margin=0.85in]{geometry}
\usepackage{array,booktabs,longtable,enumitem}
\usepackage[hidelinks]{hyperref}
\newcolumntype{L}[1]{>{\raggedright\arraybackslash}p{#1}}
\setlength{\parindent}{0pt}
\setlength{\parskip}{5pt}
\setlength{\emergencystretch}{3em}
\setlist{nosep,leftmargin=*}
\newcommand{\R}{\mathbb R}
\newcommand{\one}{\mathbf 1}
\newcommand{\DM}{\mathrm{DM}}
\newcommand{\MB}{\mathrm{MB}}
\newcommand{\sa}{\operatorname{sa}}
\newcommand{\lin}[1]{\textbf{L#1}}
\newcommand{\sha}[2]{\par{\small\ttfamily #1\allowbreak #2}\par}
\title{A08: Independent mathematical audit of P08\\Matched bilinear analysis and exact comparison}
\author{Independent reconstruction and adversarial review}
\date{3 October 2026}
\begin{document}
\maketitle

\section{Verdicts and required correction}
\begin{tabular}{@{}ll@{}}
\textbf{MB FOUNDATION VERDICT:}&\textbf{VERIFIED}\\
\textbf{MB SPECTRAL VERDICT:}&\textbf{VERIFIED}\\
\textbf{SAME-TUPLE COMPARISON VERDICT:}&\textbf{MINOR REVISION}\\
\textbf{PARITY/OBSTRUCTION VERDICT:}&\textbf{VERIFIED}\\
\textbf{OVERALL VERDICT:}&\textbf{MINOR REVISION}
\end{tabular}

All four mathematical theorem groups survive independent reconstruction.
The exact source nevertheless contains one false explicit example label:
at the universal rational witness, the MB selected representative for
$n=8$ is $2$, not $1$. This is required correction E01. It does not
invalidate the universal opposite-stability theorem, its strictness,
the all-pair theorem, or either parity theorem. The comparison and overall
verdicts conservatively retain MINOR REVISION until this source error is
corrected in the original development chat and audit closure is recorded.

\paragraph{E01: sole required correction, P08 \lin{534--535}.}
Replace the two sentences in that range with:
\begin{quote}
For $n=5,6,7$ this witness selects MB representative $1$; for $n=8$
it selects representative $2$. Other sizes follow the exact finite-grid
rule and need not retain either label.
\end{quote}
At $n=8$ the representative real parts, in order $m=0,1,2,3,4$, are
\[
 \left(-\frac14,\ -\frac18+\frac{\sqrt2}{2},\ 1,
             -\frac18-\frac{\sqrt2}{2},\ -\frac94\right).
\]
Their unique maximum is $1$ at representative $2$, giving the Fourier
labels $\{2,6\}$. In particular the erroneous label cannot be explained
by conjugate-pair notation or a tie. No change to the witness tuple,
equilibrium amplitudes, dispersion, or theorem statement is required.

\paragraph{Historical validation provenance: reproduced separately.}
The companion was absent from the attached 98-member archive but was
located in the existing workspace as \path{output/p08/P08_validation_support.zip}.
It was authenticated against all 98 current project member hashes and
the exact input ZIP hash. After completing the independent derivations
and checker, the original checker was inspected and rerun separately.
Its counts were independently derived from the code's loops before
comparison; all 271/11/280/28/198 counts and all numerical results
reproduce exactly. This removes the provisional missing-companion
limitation considered during the audit. There is no unresolved A08
provenance issue. Earlier A04/A05 historical notes are unchanged.
Crucially, the historical checker did not assert the incorrect $n=8$
label in E01: its passing run does not make that sentence true.

No P08 proof, project manifest, registry, compendium, previous audit, or
later proof is edited or promoted by this audit. In particular an audit
verdict is not a change to P08's current DRAFT lifecycle. No literature
survey, novelty promotion, or proof after P08 is undertaken.

\section{Exact current input and authentication}
The authority is the supplied \path{cyclic_network_project_P08_draft.zip},
whose SHA-256 is
\sha{3cac0e6676575e1c652454a9e28418c04}{ac5a553d4550c6f8d27f2d873e71e44}
Its current manifest entry explicitly selects
\path{proofs/P08_matched_bilinear_analysis.tex}, \textbf{P08 v1.0,
2026-10-03}, \textbf{39,287 bytes, 846 lines}. Its full-file SHA-256 is
\sha{d21a77de84df4c40c2995d30ab7401307}{95adca6c72f52fef9453ad5ddfb18ac}
This agrees exactly with the manifest's final P08 record. The body hash,
from the first line-anchored section command through EOF, is
\sha{3c4393e8a8285f946a67634addbd71d61}{7985f7f153f4963fdcbde90ca123a0d}
All P08 source ranges in this audit count comments and blank lines in
that exact file, not the older placeholder or a reconstructed proof.

The pre-existing local project directory contained stale versions,
including a 34-line P08 placeholder. They were excluded. A separate
working extraction retained the root name \path{cyclic_network_project/}
and all 98 archive members were authenticated against the archive bytes.
The delivered audit also retains the exact manifest slug, hence its name
is \path{audits/A08_matched_bilinear_analysis.tex}, rather than introducing
a new proof slug. The audit-only ZIP contains no replacement proof file.

Before mathematical review, the current manifest, registry, frozen
theorem package, final novelty verdict, novelty decisions, dependency
map, canonical model, notation, and compendium source were located and
consulted. Current P01--P05 and A01--A05 were located, together with P08;
the shared macro file was checked for conventions. Full hashes, sizes
and line counts of these 21 inputs are in
\path{A08_support/input_hashes.json}; the complete 98-member inventory is
in \path{A08_support/archive_member_hashes.json}.

All five prerequisite full-file hashes match the current manifest.
Their body hashes match the current proof-header closure records. P01's
larger-root wording, P02's ordered branch labels, and P05's corrected
candidate-versus-selected-set sentence are present. The A01/A02/A05
original minor verdicts are preserved historical evidence, with the
required changes closed in the current proofs. P03/P04 have no required
mathematical correction. A04/A05 U01 limitations remain historical and
are not claimed resolved here. P02's historical comment saying that A02
source was absent is superseded by the current manifest/compendium record
and the actual restored source; it does not prevent authentication.

The frozen order is P01--P05, P08, P07, P09, P10, then P11--P14;
P06 is later and conditional. The mathematical comparison fixes the
same positive $(a,k,c,g,h)$, $n,q$, and time variable, at each law's own
specified positive synchronous equilibrium. The novelty files withhold
priority clearance; this audit does not revisit that separate question.

\section{Independent MB foundation}
For canonical $n\ge5$, $1\le q<n$, and $a,k,c,g,h>0$, the original MB
field is
\[
 F_i(y)=ay_i-ky_i^2+cy_i y_{i-1}-gy_i y_{i+q}-hy_i y_{i-q}
       =y_i L_i(y),
\]
where $L_i=a-ky_i+cy_{i-1}-gy_{i+q}-hy_{i-q}$.
It is polynomial on $\R^n$, hence locally Lipschitz. Standard local
existence, uniqueness and continuation give a unique maximal classical
solution through every real initial state. For every finite time within
that solution's interval,
\[
 y_i(t)=y_i(0)\exp\!\left(\int_0^t L_i(y(r))\,dr\right).
\]
The exponent is finite on every compact solution interval. Thus each
coordinate preserves its sign, zero coordinates remain zero, and
nonnegative data remain nonnegative. This proves coordinate-face
invariance even with negative other coordinates. Positivity is only on
the maximal interval; no positive asymptotic lower bound or global
continuation follows from this identity alone. P08 respects that scope.

On $y=s\one$ every component equals
\[
 as-ks^2+cs^2-gs^2-hs^2=as-Ms^2,\qquad M=k+g+h-c.
\]
Scalar uniqueness and vector uniqueness identify the lift, and their
maximal intervals coincide: extension of either supplies extension of
the other. The real synchronous equilibria are $0$ and $a/M$ for
$M\ne0$, with only $0$ for $M=0$. Since $a>0$, positive synchrony is
equivalent to $M>0$, when its sole positive amplitude is $y^*=a/M$.
For $M<0$ the other real root is negative and is not a nonnegative state.

For $s_0>0$, setting $z=1/s$ gives $z'=-az+M$ and hence
\[
 s(t)=\frac{s_0e^{at}}{1+(Ms_0/a)(e^{at}-1)}.
\]
Here and in the global dynamical conclusions, time is forward from
$t=0$. When $M>0$ this is global, bounded by
$\max(s_0,a/M)$, and tends to $a/M$. When $M=0$ it is exactly
$s_0e^{at}$: global, exponential and unbounded as $t\to\infty$,
without finite-time blow-up. When $M=-|M|<0$, the denominator first
vanishes at
\[
 t_b=\frac1a\log\!\left(1+\frac{a}{|M|s_0}\right)>0.
\]
The denominator is positive before $t_b$ and the numerator is positive
at it, so $s(t)\to+\infty$ as $t\uparrow t_b$. This is an actual
full-system solution because of invariant synchrony. Zero data remain
zero for all three signs of $M$; no assertion about every nonsynchronous
trajectory when $M\le0$ is warranted or made.

\subsection{Maximum estimate, ties, and exact implications}
For $Y=\max_i y_i$ and $I(t)=\{i:y_i=Y\}$, a finite first-order Taylor
expansion at a fixed time excludes the nonmaximal coordinates locally
and yields
\[
 D^+Y=\max_{i\in I(t)}\dot y_i.
\]
For every such $i$, $0\le y_{i-1},y_{i+q},y_{i-q}\le Y$, giving
\[
 \dot y_i\le aY-kY^2+cY^2=aY-(k-c)Y^2.
\]
This proof covers multiple maximizing coordinates and every allowed
offset collision; it bounds each original term and never overwrites a
coincident contribution.

If $k>c$, logistic comparison gives $Y(t)\le\max(Y(0),R)$,
$R=a/(k-c)$, global continuation, invariance of $[0,R]^n$, uniform
attraction of bounded sets, and finite absorption into every
$[0,R+\varepsilon]^n$. One may dominate every initial maximum from a
bounded set by the same scalar solution starting at its supremum or
at $R$, whichever is larger. This avoids needing a positive lower bound
on that supremum. The estimate alone does not show finite absorption
into the exact comparison cube.

If $k=c$, it gives $Y(t)\le Y(0)e^{at}$ and thus global forward
existence, but no bound uniform over $t\ge0$. If $k<c$, its scalar
upper comparison can blow up, so this estimate alone yields neither
global continuation nor failure thereof. It is not a necessary
condition for boundedness. The synchronous examples do prove that
$M>0$ is necessary for \emph{every} nonnegative trajectory to be bounded,
and $M\ge0$ is necessary for \emph{every} such trajectory to be global.
Neither converse is proved or asserted. These are exactly P08's scopes.

\section{Independent Jacobian and full Fourier spectrum}
Differentiate before imposing any equilibrium relation. With cyclic
Kronecker deltas, the derivative is
\begin{align*}
 \frac{\partial F_i}{\partial y_j}
 ={}&(a-2ky_i+cy_{i-1}-gy_{i+q}-hy_{i-q})\delta_{ji}\\
 &+cy_i\delta_{j,i-1}-gy_i\delta_{j,i+q}-hy_i\delta_{j,i-q}.
\end{align*}
None of the three neighbor offsets is zero in the canonical range.
At $y=s\one$ its diagonal is
$a-2ks+cs-gs-hs=a-(k+M)s$, and the three neighbor contributions
are $cs,-gs,-hs$. With $(Pz)_i=z_{i-1}$ this gives
\[
 DF(s\one)=[a-(k+M)s]I+s(cP-gP^{-q}-hP^q).
\]
Only at $s=y^*=a/M>0$ may one substitute $a=My^*$ to obtain
\[
 J_{\MB}=y^*(-kI+cP-gP^{-q}-hP^q).
\]

The verified convention is $v^{(m)}_j=n^{-1/2}e^{\mathrm i\theta_m j}$,
$\theta_m=2\pi m/n$. Direct application gives
$P^\ell v^{(m)}=e^{-\mathrm i\ell\theta_m}v^{(m)}$. Therefore
\begin{align*}
 \lambda_m^{\MB}&=y^*[-k+ce^{-\mathrm i\theta_m}
                         -ge^{\mathrm i q\theta_m}-he^{-\mathrm i q\theta_m}],\\
 \Re\lambda_m^{\MB}&=y^*[-k+c\cos\theta_m-(g+h)\cos(q\theta_m)],\\
 \Im\lambda_m^{\MB}&=y^*[-c\sin\theta_m+(h-g)\sin(q\theta_m)].
\end{align*}
The signs of the imaginary contributions are respectively $-c,-g,+h$.
The geometric sum makes the Fourier matrix unitary, so these $n$ labels
give the complete eigenvalue list with algebraic multiplicity. A unitary
diagonalization proves normality and semisimplicity even if several
labels give the same complex eigenvalue.

At $m=0$ the matrix formula gives
\[
 \lambda_0^{\MB}=y^*(-k+c-g-h)=-My^*=-a.
\]
Independently, $f'(s)=a-2Ms$ for the scalar restriction, so
$f'(a/M)=-a$. The tangent direction is $\one$, identifying exactly
the uniform mode. Consequently positive MB synchrony can be linearly
unstable only in a nonuniform direction. This does not classify a
steady or oscillatory bifurcation, or imply nonzero frequency.

\subsection{All collision cases}
Contributions add in the original derivative, giving
\[
\begin{array}{c|l|l}
q&J_{\MB}/y^*&\Im\lambda_m^{\MB}/y^*\\\hline
1&-kI+(c-h)P-gP^{-1}&(-c+h-g)\sin\theta_m\\
n-1&-kI+(c-g)P-hP^{-1}&(-c+g-h)\sin\theta_m\\
n/2&-kI+cP-(g+h)P^{n/2}&-c\sin\theta_m
\end{array}
\]
The final row requires even $n$ and uses $P^{-n/2}=P^{n/2}$;
its restraint real part is $-(g+h)(-1)^m$. For even $n$, the Nyquist
value is $y^*[-k-c-(g+h)(-1)^q]$. All zero-mode values remain $-a$.
The general formula is valid without coprimality; the collision checks
do not impose extra restrictions on the $q=2$ comparisons.

\subsection{Exact stability and the additional global bound}
Put $S=g+h$ and
\[
 K_{n,q}=\max_m[c\cos\theta_m-S\cos(q\theta_m)].
\]
At positive synchrony $\sa(J_{\MB})=y^*(K_{n,q}-k)$. Strict stability
is exactly $k>K_{n,q}$ and instability is exactly $k<K_{n,q}$.
At equality every real part is nonpositive with a zero maximum.
Normality gives $\|e^{tJ}\|_2=e^{t\sa(J)}$, so the linear system is
Lyapunov stable but not asymptotically stable. Nonlinear boundary
stability remains undecided. A nonzero spectral maximum alone is not
equivalent to hyperbolicity because another real part could be zero.

For the stronger full-flow estimate, let $B=-kI+cP-gP^{-q}-hP^q$,
$\delta=k-K_{n,q}>0$, and $V=\sum_i y_i$. Its symmetric part has
eigenvalues $-k+c\cos\theta_m-S\cos(q\theta_m)\le-\delta$; its
skew part contributes zero to $y^{\mathsf T}By$. Hence
\[
 V'=aV+y^{\mathsf T}By\le aV-\delta\|y\|_2^2
                        \le aV-\frac\delta n V^2.
\]
Nonnegativity bounds every coordinate by $V$, so this excludes finite
escape and yields a bounded absorbing simplex at any radius larger than
$na/\delta$. The simplex at the exact radius is invariant and attracts
bounded sets. Moreover $K_{n,q}\ge c-S$ from mode zero, so $M>0$ follows.
Thus strict MB spectral stability is a sufficient global dissipativity
condition for the nonnegative full flow. It is not being asserted
necessary for full-flow boundedness.

The witness $(9/2,7,8,1,1)$ has $M=1$, $y^*=9/2$, $x^*=1$,
$\mu=23/2>0$, and zero original residuals. For $q=2$,
$K_{n,2}=\max_u(2+8u-4u^2)=6$ on every grid, uniquely at $u=1$.
Thus MB is stable and dissipative despite $k=7<c=8$. Its maximum is
$-9/2$. DM has $d=-13,C=4,G=1/4,H=1$, vertex $4/3$, and maximum
$-39/4$ at zero. This verifies the full spectra via the formulas above,
not only a selected-mode label. Conversely $k>c$ alone need not imply
MB spectral stability: for example $n=8,q=2,(a,k,c,g,h)=(1,2,1,2,2)$
has $M=5>0$ and the $u=0$ MB mode has normalized real part $2>0$.
This auxiliary audit example simply checks P08's scope statement.

\section{Exact step-two grid selection and stability}
From now on the comparison is $q=2$. Let
$u_m=\cos\theta_m$, $N=\lfloor n/2\rfloor$ and
$1=u_0>u_1>\cdots>u_N$. Direct substitution gives
\begin{align*}
 p(u)&=\Re\lambda^{\MB}/y^*=-k+S+cu-2Su^2\\
     &=-k+S+\frac{c^2}{8S}-2S\left(u-\frac{c}{4S}\right)^2.
\end{align*}
It is strictly concave since $S>0$. The unconstrained vertex is
$v=c/(4S)>0$ and the constrained continuous maximizer on $[-1,1]$
is $\min(1,v)$. Both positivity and strict concavity needed by P05's
generic grid result have thus been established anew for MB.

Define $\beta_j=(u_j+u_{j+1})/2$. These midpoints strictly decrease,
and the cosine addition formula gives
\[
 \beta_j=\cos(\pi/n)\cos((2j+1)\pi/n).
\]
Exactly the midpoints $0\le j<J$, $J=\lfloor(n+1)/4\rfloor$, are
positive. To see the count, $\cos(\pi/n)>0$ and the other angle lies
between $0$ and $\pi$, so positivity is precisely $2j+1<n/2$.
Only these switches are reachable with $v>0$.

The nearest-cosine rule follows directly from the completed square,
with closed cells
\[
 [\beta_0,\infty),\quad [\beta_j,\beta_{j-1}]\ (1\le j<N),
 \quad(-\infty,\beta_{N-1}],
\]
intersected with $v>0$. Equality at an attainable midpoint gives only
the adjacent representatives; a third or nonadjacent maximal pair
would contradict strict concavity at an intermediate grid point.
Putting $\chi_{\MB}=4S/c=1/v$, $\alpha_j=1/\beta_j$ for $j<J$,
the strict selections are
\[
\begin{array}{c|c}
0<\chi_{\MB}<\alpha_0&0\\
\alpha_{j-1}<\chi_{\MB}<\alpha_j, 1\le j<J&j\\
\chi_{\MB}>\alpha_{J-1}&J
\end{array}
\]
and $\chi_{\MB}=\alpha_j$ selects $\{j,j+1\}$. An independent
subtraction confirms orientation:
\[
 p(u_j)-p(u_{j+1})=c(u_j-u_{j+1})(1-\chi_{\MB}\beta_j).
\]
The middle row is empty at $J=1$. The endpoint $v=1$ itself is not
a discrete tie. The exact reachable set is $0,\ldots,J$; the Nyquist
singleton cannot be selected. A nonzero representative represents two
conjugate labels; a tie involving zero gives three labels, another
reachable tie four. These counts are not eigenvalue simplicity claims.

MB realizability is independent of DM realizability: given any $v>0$,
take $g=h=1/2$, $c=4v$, $k>\max(0,c-1)$ and any $a>0$.
Then $M=k+1-c>0$ and the desired vertex is realized. This covers every
strict cell and every positive-midpoint tie.

The four terminal grid cases are
\[
\begin{array}{c|c|c|c}
n&J&u_J&\text{nearest at formal }v=0\\\hline
4\ell&\ell&0&\ell\\
4\ell+1&\ell&\sin(\pi/(2n))&\ell\\
4\ell+2&\ell&\sin(\pi/n)&\ell,\ell+1\\
4\ell+3&\ell+1&-\sin(\pi/(2n))&\ell+1
\end{array}
\]
The negative reachable terminal cosine is unique to the last class.
For $n=4\ell+2$, the formal midpoint $0$ is unattainable at finite
positive coefficients; the positive and negative nearest cosines differ
in MB real growth by $2y^*c\sin(\pi/n)>0$. Dividing by $y^*$ removes
the amplitude dependence; the unnormalized gap has no uniform
positive lower bound along arbitrary parameter limits.

Because $y^*>0$, stability in a selected cell $j$ is precisely
$k>S+cu_j-2Su_j^2$. Either representative at a tie gives the same
bound. Reverse inequality is instability and equality is zero abscissa.
If zero is among the maxima, $M>0$ already gives maximum $-a<0$.
All equality boundaries and signs in P08's selection theorem pass.

\section{Common equilibrium domain and structural distinctions}
The original DM synchronous field, after division by a positive amplitude,
is
\[
 B(x)=a-(k+h)x+\frac{c-g}{1+x}.
\]
With $A=k+h$, $\mu=a+c-g$, $b=a-A$, and $D=b^2+4A\mu$,
the nonzero equilibrium polynomial is $Q(x)=Ax^2-bx-\mu$.
P02's complete classification gives the common positive-root domain
\[
 M>0,\qquad
 \mu>0\quad\text{or}\quad(\mu=0,b>0)
 \quad\text{or}\quad(\mu<0,b>0,D\ge0).
\]
It includes a boundary and a double-root locus and is not wholly open.
Every positive branch must be specified when two roots occur. On
$\mathcal C_+=\{(a,k,c,g,h)>0:M>0,\mu>0\}$, the sole positive
DM root is $x_+=(b+\sqrt D)/(2A)$, and it is simple. P08 correctly
uses this domain for the main all-pair and unique-branch contrasts and
names $x_-$ for the distinct lower-branch construction.

At a specified positive root let $s=1+x$. Independent differentiation
and then use of $B(x)=0$ give
\[
 d=-kx-\frac{c(1+2x)}{s^2},\quad C=\frac cs,\quad
 G=\frac{gx}{s^2},\quad H=hx,
\]
\[
 \lambda_m^{\DM}=d+Ce^{-\mathrm i\theta_m}
                       +Ge^{2\mathrm i\theta_m}-He^{-2\mathrm i\theta_m}.
\]
Thus $r_m^{\DM}=d+Cu_m+(G-H)(2u_m^2-1)$, with imaginary part
$-C\sin\theta_m+(G+H)\sin(2\theta_m)$. The selection controls are
\[
 \chi_{\DM}=\frac{4(H-G)}C
 =\frac{4x(hs^2-g)}{cs},\qquad
 \chi_{\MB}=\frac{4(g+h)}c.
\]
In particular
\[
 \frac{\chi_{\DM}}{\chi_{\MB}}
 =\frac{x(hs^2-g)}{s(g+h)}.
\]
There is no universal ordering of them. At fixed $k,c,g,h$, MB's
control is independent of $a$; DM's can move through $x$. Arbitrary
positive pairs can in fact be realized, as proved in the next section.

DM's curvature is $4(G-H)$: convex, linear, concave for respectively
$h<g/s^2$, $h=g/s^2$, $h>g/s^2$. MB's is always
$-4y^*(g+h)<0$. In the nonconcave DM cases,
\[
 r(1)-r(u)=(1-u)[C+2(G-H)(1+u)]>0\quad(-1\le u<1),
\]
so zero is uniquely dominant. The common-domain tuples
$x=y=1,c=g=4,k=1,a=1+h$, $h=1/2,1,2$, have $M=\mu=1+h>0$
and $G-H=1-h$, proving all three DM curvature signs are attainable.

DM has $F_i|_{x_i=0}=cx_{i-1}$, whereas MB's same coordinate face is
invariant. P01's directed predecessor propagation makes every coordinate
strictly positive for $t>0$ from a nonzero nonnegative DM initial state.
The DM uniform eigenvalue is $xB_x=\mp x\sqrt D/(1+x)$ on the
upper/lower distinct roots and is zero at a positive double root;
MB's is always $-a$. No equality of the two amplitudes or derivatives
has been assumed.

P01 supplies global nonnegative existence and bounded absorption for
every canonical DM tuple. MB has the sufficient dissipative regimes
already proved, but the two synchronous counterexamples prevent a
whole-class claim. P08's phrase ``structurally dissipative'' is explicitly
qualified by ``over the entire canonical class,'' with its intended
meaning supplied by these preceding mathematical statements. It does
not assert that every MB tuple is unbounded.

Matching controls the node number, directed offsets, coefficient tuple
and count, coefficient roles and original time variable. It does not
calibrate amplitudes, effective derivatives, nonlinear effects or empirical
behavior. In this $q=2$ comparison the successor contribution has
opposite derivative signs in the two laws. The fixed-source limit
$cz/(1+r)\to cz$ versus $crz\to0$ as $r\downarrow0$ also passes.
If this sign discussion is quoted for collision offsets, it must refer
to the individual contribution; net coincident entries are sums.
No biological equivalence or quantitative superiority is claimed.

\section{Every ordered pair with both models stable}
The precise quantifier is: for every fixed $n\ge5,q=2$ and every
$(j_{\DM},j_{\MB})\in\{0,\ldots,J\}^2$, a nonempty open subset
of physical five-dimensional $\mathcal C_+$ has that unique pair of
representatives and both equilibria Hurwitz. This includes zero and
terminal representatives in all parity classes, not only interior modes.

Choose strictly interior positive controls $\eta,\zeta$ in the desired
DM and MB cells, using $\alpha_0/2$, the arithmetic means of adjacent
$\alpha$'s, and $2\alpha_{J-1}$ for the two endpoints and interior
cells as appropriate. Set
\[
 c=4,\quad S=\zeta,\quad x=1+\eta/\zeta,\quad s=1+x,
 \quad h=\frac{\zeta+s\eta/x}{s^2+1},\quad g=\zeta-h.
\]
Certainly $h>0$; moreover $h<\zeta$ is equivalent to
$\eta<\zeta xs$, which holds since $xs>\eta/\zeta$. Thus $g,h>0$.
Choose
\[
 k>\max\{c+S,(c-g)/(xs)-h\},\qquad
 a=(k+h)x-(c-g)/s.
\]
Then $a>0$, $M>0$, and
$\mu=x[k+h+(c-g)/s]>0$, because $k>S>g$ and $s>1$.
Substitution in the \emph{original} fields gives
\[
 F_i^{\DM}(x\one)=x[a-(k+h)x+(c-g)/s]=0,
 \qquad F_i^{\MB}((a/M)\one)=(a/M)[a-M(a/M)]=0.
\]
Now $(s^2+1)h-\zeta=s\eta/x$, so
$x(hs^2-g)/s=\eta$. With $c=4$ this proves exactly
$\chi_{\DM}=\eta$ and $\chi_{\MB}=\zeta$.

All MB real parts divided by $y^*$ are at most $-k+c+S<0$.
For every DM mode, using $u\le1$ and $|\cos2\theta|\le1$,
\[
 \frac{r_m^{\DM}}x
 \le-k-\frac c{s^2}+\frac g{s^2}+h<-k+S<0.
\]
The $c$ bound follows from
$-(1+2x)+su\le-x$. The latter strict inequality is valid for every
positive $x,c,g,h$. The large-$k$ operation adjusts $a$ while retaining
the prescribed $x,c,g,h$, so neither selection control moves. This
addresses the principal possible failure in the all-pair construction.

Every chosen control lies strictly inside its cell; no endpoint or
parity exception is omitted. Nonempty openness in all five physical
parameters follows from scalar branch regularity, proved below. Distinct
nonzero representatives first become possible at $J=2$, namely $n=7$.

At $n=8$ the separate rational tuple $(113/4,9,8,1,1)$ gives
$M=3$, $x=3$, $y=113/12$, $\mu=141/4>0$. Both original divided
residuals vanish:
\[
 113/4-10\cdot3+7/4=0,\qquad113/4-3(113/12)=0.
\]
DM has $(d,C,G,H)=(-61/2,2,3/16,3)$ and
$\chi_{\DM}=45/8>2\sqrt2=\alpha_1$, selecting $2$ with
maximum $-443/16$. MB has $\chi_{\MB}=1$, selecting $0$ with
maximum $-113/4$. The full spectral formulas certify both all-mode
stability assertions, including unequal own equilibrium amplitudes.

\section{Universal forward stability contrast and E01}
Take $(a,k,c,g,h)=(1/4,1/8,1,1,1/8)$ for any $n\ge5,q=2$.
Here $M=\mu=1/4>0$. The original synchronous residuals at amplitude
one are respectively
\[
 \frac14-\frac18+\frac12-\frac12-\frac18=0,
 \qquad\frac14-\frac18+1-1-\frac18=0.
\]
Thus $x=y=1$ by independent calculation and uniqueness, rather than
an imposed matching assumption. For DM,
\[
 (d,C,G,H)=(-7/8,1/2,1/4,1/8),\qquad
 r^{\DM}(u)=-1+u/2+u^2/4.
\]
Its gap to the uniform value is
\[
 r^{\DM}(1)-r^{\DM}(u)=(1-u)(3+u)/4>0\quad(u<1),
\]
and $r^{\DM}(1)=-1/4$. This proves the exact spectral abscissa
and unique representative for every size, including arbitrarily large $n$.

MB has $S=9/8$, $y=1$ and
\[
 p(u)=1+u-\frac94u^2,\qquad v_{\MB}=2/9.
\]
An actual grid label $J=\lfloor(n+1)/4\rfloor$ satisfies
$-1/4<u_J\le1/2$. Indeed the four parity formulas apply; the only
negative case has $n\ge7$ and
$|u_J|=\sin(\pi/(2n))<\pi/14<1/4$. In the $4\ell+2$ case
the largest value is $\sin(\pi/6)=1/2$ at $n=6$, and the other
nonnegative cases are smaller. Consequently
\[
 p(u_J)\ge1-\frac14-\frac94\left(\frac12\right)^2
          =\frac3{16}>0.
\]
This is a conservative valid uniform lower bound, even though its
separate worst bounds on $u$ and $u^2$ do not occur simultaneously.
Since $J>0$, it establishes
$\max_{m\ne0}\Re\lambda_m^{\MB}\ge3/16$ without a continuous
grid approximation or finite sample as proof.

There is no representative tie for any $n$. A tie would require
$\beta_j=2/9$. But
$4\beta_j=2\cos(2\pi j/n)+2\cos(2\pi(j+1)/n)$ is a sum of
roots of unity and their inverses, hence an algebraic integer. It
would equal the noninteger rational $8/9$, impossible: the rational
root theorem says a rational root of a monic integer polynomial is an
integer. This establishes strict finite-grid uniqueness. The selected
representative is nonzero because mode zero equals $-a=-1/4$.

At the witness $B_x=-1/4\ne0$ and the MB scalar derivative is
$-1/4\ne0$. Positivity, $M>0$, $\mu>0$, both stability margins, and
both representative gaps are strict for each fixed $n$. The physical
open-neighborhood claim therefore passes. The draft explicitly allows
this neighborhood to depend on $n$; it does not prove or claim one
uniform neighborhood preserving the selected representatives for all $n$.

\paragraph{The incorrect finite-size sentence.}
For $n=5$ the two positive-vertex candidates are $0,1$ and instability
excludes zero. For $n=6$, $p(1/2)=15/16$ is maximal at $1$. For
$n=7$, comparing the two reachable nonzero representatives gives
\[
 p(u_1)-p(u_2)=(u_1-u_2)\left[1-\frac94(u_1+u_2)\right]>0,
\]
because the rational cosine intervals verified below imply
$u_1+u_2<41/100<4/9$. Thus $1$ is correct for $n=5,6,7$.
For $n=8$, $u_2=0$ and $p(0)=1$ exceeds every other representative
value listed in E01. Its strict next-mode gap is
$9/8-\sqrt2/2>0$. P08 \lin{534} is therefore false for $n=8$,
even though all surrounding all-$n$ theorem conclusions are true.

\section{Reverse stability on the unique positive DM domain}
For any positive DM equilibrium put $s=1+x$, $u=\cos\theta_m$,
$w=-\cos(2\theta_m)$. Independent substitution gives
\begin{align*}
 \frac{r_m^{\DM}}x&=-k+(h-g/s^2)w-
                  \frac{c(1+2x-su)}{xs^2},\\
 \frac{r_m^{\MB}}{y}&=-k+(g+h)w+cu.
\end{align*}
Here $1+2x-su\ge x>0$ for $u\le1$.

\subsection{Necessity and equality cases}
On $\mathcal C_+$ the uniform DM value is strictly negative. If
$G-H\ge0$, zero uniquely maximizes, proving DM stable directly.
If $G-H<0$, then $h-g/s^2>0$. Any nonnegative DM mode must have
$w>0$, since both $-k$ and the final term are strictly negative.
When also $u\ge0$, subtraction gives
\[
 \frac{r_m^{\MB}}y-\frac{r_m^{\DM}}x
 =g(1+s^{-2})w+cu+\frac{c(1+2x-su)}{xs^2}>0.
\]
Hence even a \emph{zero} DM mode forces a positive MB mode. For
$n\equiv0,1,2\pmod4$, every reachable DM representative has a
nonnegative cosine. Applying this to a maximizing representative
excludes any nonnegative DM spectral abscissa if MB is strictly stable.
Ties among DM representatives cause no loophole: choose any maximizer;
every reachable cosine in these classes is nonnegative. Equality
$u=0$ still gives a strictly positive gap. The case $G-H=0$ was
already handled by unique uniform selection. Thus the stronger strict
implication in the draft is correct, not merely an absence of a sampled
unstable point.

For $n\equiv3\pmod4$, the terminal representative has negative
cosine. This is the exact obstruction-proof feature, rather than a
parity-table guess. The gap argument cannot force MB positivity in
that mode, and the following construction realizes the exception.

\subsection{Sufficiency for every size in the remaining class}
Let $n\equiv3\pmod4$, $n\ge7$,
$\tau=\sin(\pi/(2n))\in(0,1/4)$,
$W=1-2\tau^2>7/8$. Then $u_J=-\tau$,
$u_{J-1}=3\tau-4\tau^3$, and $\beta_{J-1}=\tau W$.
Choose
\[
 h=1,\quad c=\tau,\quad g=\tau^2/4,\quad
 k=1-5\tau^2/2,\quad x=2/\tau,\quad
 a=(k+1)x-(\tau-\tau^2/4)/(1+x).
\]
All coefficients are positive: $k>27/32$, $x>8$, and the positive
large term in $a$ exceeds $8$ while the subtracted term is below $1/4$.
Also $M=2-\tau-9\tau^2/4>0$, and $\mu>0$ because $c>g$ and
$a>0$. The original DM residual is zero by the definition of $a$;
the MB root is its own positive $y=a/M$.

For MB,
$v_{\MB}=\tau/(4+\tau^2)<\tau/4<\tau W=\beta_{J-1}$,
so $J$ is uniquely selected. Its normalized maximum is exactly
\[
 -k+(1+\tau^2/4)W-\tau^2
 =-\tau^2(1/4+\tau^2/2)<0.
\]
For DM, the same label gives
\[
 \frac{r_J^{\DM}}x=\frac{\tau^2}{2}-\frac{gW}{s^2}
                  -\frac{\tau(1+2x+s\tau)}{xs^2}.
\]
Because $1+2x+s\tau\le(2+\tau)s$, $s>x$, and $x=2/\tau$,
the sum of the two subtractions is less than
\[
 \frac{g+\tau(2+\tau)}{x^2}
 <\frac{1/64+9/16}{4}\tau^2
 =\frac{37}{256}\tau^2.
\]
Thus $r_J^{\DM}/x>91\tau^2/256>0$. Concavity holds, and
\[
 v_{\DM}=\frac{\tau}{4x(s-g/s)}<\frac{\tau}{256}
                         <\beta_{J-1},
\]
since $x>8$ and $s-g/s>8$. Therefore DM also uniquely selects $J$.
Every inequality is strict for every such finite $n$, although some
margins approach zero as $n$ grows. Branch simplicity and finite-mode
continuity imply a physical open region for each fixed size. This
proves sufficiency and completes the stated iff theorem.

\subsection{Exact rational seven-node witness}
For $(a,k,c,g,h)=(20301/1100,17/20,1/2,1/100,1)$,
$M=34/25$, $x=10$, $y=20301/1496$, and $\mu>0$.
The divided original residuals are exactly
\[
 20301/1100-(37/20)10+(49/100)/11=0,
 \quad20301/1100-(34/25)(20301/1496)=0.
\]
The three ordered nontrivial cosines are all the roots of
$8u^3+4u^2-4u-1$; this follows by applying the seventh-root geometric
sum and Chebyshev identities. Exact rational evaluation gives signs
$(-,+)$ on $(31/50,63/100)$, $(+,-)$ on
$(-23/100,-11/50)$, and $(-,+)$ on $(-91/100,-9/10)$.
Three disjoint sign changes exhaust a cubic's degree, so the intervals
isolate the roots and cosine ordering assigns $u_1,u_2,u_3$.

The normalized MB real parts for representatives $1,2,3$ are
\[
 -17/20+u_1/2-101u_2/100,\quad
 -17/20+u_2/2-101u_3/100,\quad
 -17/20+u_3/2-101u_1/100.
\]
Their strict rational upper bounds are respectively
$-3027/10000$, $-409/10000$, $-19262/10000$; the zero mode is
$-34/25$. All seven eigenvalues therefore have negative real part.
DM's coefficients are
$(d,C,G,H)=(-1039/121,1/22,1/1210,10)$, and
\[
 \frac{r_2^{\DM}}{10}
 >-\frac{17}{20}-\frac{21}{2420}-\frac{23}{22000}
       +\frac{12099}{12100}\frac9{10}
 =\frac{9729}{242000}>\frac1{25}.
\]
Both vertices, $25/202$ and $55/48396$, lie below
$39/200<\beta_1$ and above the negative $\beta_2$, so both selected
representatives are uniquely $2$. The full imaginary parts follow by
substituting these exact coefficients in the displayed full spectra;
the accompanying results list all seven complex eigenvalues. No decimal
cosine or numerical sign is needed for the proof.

\section{Full branch-labelled domain and stronger obstruction}
If $4\mid n$, the MB mode $n/4$ has $u=0$ and normalized real part
$-k+S$. Strict MB stability forces $k>S$. At \emph{any} positive DM
equilibrium, on either allowable branch, the bound already derived gives
\[
 r_m^{\DM}/x\le-k-c/s^2+g/s^2+h<-k+S<0.
\]
This used only $x>0$ and the equilibrium identity, not $\mu>0$.
It covers the positive upper root, a putative positive lower root and a
putative positive double root. Indeed the latter two cannot coexist
with these hypotheses because their uniform eigenvalues are positive
or zero, contradicting the bound. MB still has its sole positive root
$a/M$. This proves the stronger full-domain obstruction without an
omitted branch or an equality loophole. It is a strict-stability theorem;
it does not silently make a boundary assertion when $k=S$.

If $4\nmid n$, define $\rho=\max_m[-\cos(2\theta_m)]$.
It satisfies $0<\rho<1$. Equality $\rho=1$ would require
$4m=n\pmod{2n}$ and hence $4\mid n$. For positivity, a nearest
integer to $n/4$ gives a double angle within $2\pi/n<\pi/2$ of
$\pi$ for $n\ge5$, so its cosine is negative. Let
\[
 e=(1-\rho)/24,\quad k=1,\quad c=h=e,\quad
 g=(1+e)^3+2e,\quad
 a=(1+e)(1+2e)+e/(1+e).
\]
All coefficients are positive and $M=1+g>0$. Writing $A=1+e$,
direct expansion of the original DM equation gives exactly the two roots
\[
 x_-=e,\qquad x_+=e+e/(1+e)^2,
\]
with $\mu=-Ae^2-e^2/A<0$, $b=A(x_-+x_+)>0$ and
$D=A^2(x_+-x_-)^2>0$. Thus these are the distinct positive roots
of the original synchronous equation, not only formal roots.
The branch being compared is explicitly the lower one.

At $x_-=e$, independent calculation gives
\[
 \lambda_0^{\DM}=e^2/(1+e)^2>0,\qquad
 G-H=e\left(1+2e/(1+e)^2\right)>0.
\]
Hence zero uniquely dominates and DM is unstable. For MB,
$S=(1+e)^3+3e<1+10e$ since $0<e<1$, and every real part divided
by $y$ is bounded by
\[
 -1+e+S\rho<-1+e+(1+10e)\rho
 =-(1-\rho)+e(1+10\rho)
 <-\frac{13}{24}(1-\rho)<0.
\]
This proves full MB stability with its own amplitude $y=a/M$ and
strict all-mode margin. Both DM roots are simple, so the named lower
root and inequalities persist locally. At $n=6$, $\rho=1/2$ and
$e=1/48$, giving a fully rational instance with
\[
 (a,k,c,g,h)=\left(\frac{61177}{56448},1,\frac1{48},
                       \frac{122257}{110592},\frac1{48}\right),
 \quad x_-=\frac1{48},\quad x_+=\frac1{48}+\frac{48}{2401}.
\]
The exact fractions are independently checked in the support results.

Consequently on the full branch-labelled common domain reverse
stability exists exactly when $4\nmid n$. On $\mathcal C_+$ it
exists exactly when $n\equiv3\pmod4$. For $n\equiv1,2\pmod4$
there is no unique-domain reverse contrast, but there is the proved
lower-branch contrast. P08 does not claim a complete upper-branch
classification on the remaining part of the common domain; none is
inferred here. There is no conflict between these differently scoped
iff statements.

\section{Remaining explicit examples and their full spectra}
For any $q=2$ witness, listing $x,y,d,C,G,H$ and its original tuple
determines every complex eigenvalue exactly through
\[
 \lambda_m^{\DM}=d+C\cos\theta_m+(G-H)\cos2\theta_m
        +\mathrm i[-C\sin\theta_m+(G+H)\sin2\theta_m],
\]
\[
 \lambda_m^{\MB}=y[-k+c\cos\theta_m-S\cos2\theta_m
        +\mathrm i(-c\sin\theta_m+(h-g)\sin2\theta_m)].
\]
These are full complex lists, including uniform, conjugate and Nyquist
labels; imaginary parts do not determine stability by themselves.
The support results additionally record their numerical evaluation for
every explicit common-domain rational example.

\subsection{Both unstable and different modes, only the claimed n=8 case}
Take $(16351/1100,1/2,3/2,1/100,1)$, $n=8,q=2$.
All coefficients are positive, $M=1/100$, $\mu>0$,
$x=10$, and $y=16351/11$. Both original divided residuals vanish:
\[
 16351/1100-15+(149/100)/11=0,
 \quad 16351/1100-(1/100)(16351/11)=0.
\]
DM has $(d,C,G,H)=(-1273/242,3/22,1/1210,10)$ and vertex
$55/16132<\sqrt2/4$. Its selected representative is uniquely $2$,
with
\[
 \sa(J_{\DM})/10=-\frac12-\frac{63}{2420}
                   +\frac{12099}{12100}=\frac{2867}{6050}>0.
\]
MB's vertex is $75/202$, strictly between $\sqrt2/4$ and
$(2+\sqrt2)/4$, so it uniquely selects representative $1$ and
\[
 \sa(J_{\MB})/y=-1/2+3\sqrt2/4>0.
\]
For extra explicit full-spectrum control, the representative MB real
parts divided by $y$ are
\[
 \left(-\frac1{100},-\frac12+\frac{3\sqrt2}{4},
                  \frac{51}{100},-\frac12-\frac{3\sqrt2}{4},
                  -\frac{301}{100}\right).
\]
The positive selected value exceeds $51/100$ strictly, equivalent to
$75/202>\sqrt2/4$. All DM and MB imaginary parts are given by the
preceding exact full-spectrum formulas. There is no tie ambiguity.
The small positive $M$ explains the large MB amplitude; it is never
silently identified with $x=10$. This proves an open example at $n=8$,
not an all-size theorem about two unstable models.

\subsection{Curvature witnesses, including behavior not claimed by P08}
For the three tuples $(1+h,1,4,4,h)$, $h=1/2,1,2$, both amplitudes
are $1$, $M=\mu=1+h>0$, both original residuals vanish, and
$(d,C,G,H)=(-4,2,1,h)$. Thus for all labels
\[
 r_m^{\DM}=-4+2u_m+(1-h)\cos2\theta_m,
 \quad\Im\lambda_m^{\DM}=-2\sin\theta_m+(1+h)\sin2\theta_m,
\]
\[
 r_m^{\MB}=-1+4u_m-(4+h)\cos2\theta_m,
 \quad\Im\lambda_m^{\MB}=-4\sin\theta_m+(h-4)\sin2\theta_m.
\]
At $n=8$ the DM selections are $0,0,1$, with maxima
$-3/2,-2,-4+\sqrt2$; MB selects $2$ with maxima $7/2,4,5$.
These additional checked labels simply document the examples' full
behavior; P08 uses them to prove curvature realizability and does not
claim these particular labels for every size.

\subsection{Structural growth examples outside the common domain}
The tuple $(1,1,3,1,1)$ has $M=0$, no positive MB equilibrium,
and at synchronous $s_0=1$ the full MB trajectory is $s=e^t$.
DM's positive root is $x=1$, with original residual zero and
$(d,C,G,H)=(-13/4,3/2,1/4,1)$. At $q=2$ its dispersion has vertex
$1/2$ and is strictly negative for every $u\in[-1,1]$, since its
continuous maximum is $-17/8$. This stability observation is not
needed for the global boundedness contrast.

The tuple $(1,1,4,1,1)$ has $M=-1$, again no positive MB equilibrium,
and $s=e^t/(2-e^t)$ blows up at $\log2$ from $s_0=1$.
The DM root is $x=(\sqrt{33}-1)/4>0$, since
$2x^2+x-4=0$, equivalently $1-2x+3/(1+x)=0$. Its spectral
coefficients are $d=-x-4(1+2x)/(1+x)^2$,
$C=4/(1+x)$, $G=x/(1+x)^2$, $H=x$. The master spectrum provides
every label; no MB positive-root spectrum can or should be attached.
In both examples DM is globally bounded by P01, regardless of local
spectral behavior. The MB origin has zero residual and derivative
$aI=I$, while each stated synchronous trajectory has exactly the
original full MB vector derivative. These examples are not common-root
stability comparisons.

\section{Openness, branch regularity, and logical scope}
For any positive root, $B(x)=-Q(x)/(1+x)$. Differentiating at a
root gives
\[
 B_x=-Q'(x)/(1+x)=\mp\sqrt D/(1+x_\pm).
\]
On $\mathcal C_+$, $D=b^2+4A\mu>0$, so the positive upper root
is simple and smoothly determined by all five physical parameters.
In the lower family $\mu<0,b>0,D>0$, both roots are distinct and
positive; the selected lower root has $B_x=+\sqrt D/(1+x_-)>0$.
The MB root $a/M$ is smooth where $M>0$ and its scalar derivative
is $-a\ne0$.

Every constructed strict contrast therefore continues on a locally
unique positive scalar branch under arbitrary sufficiently small
perturbations of $(a,k,c,g,h)\in(0,\infty)^5$. Strict parameter,
existence and branch conditions remain strict. For fixed $n$ there are
only finitely many Fourier values and representative comparisons, all
continuous on the branch. Their strict stability signs and strict
gaps persist on an intersection of open neighborhoods. A conjugate pair
remains tied, as required; uniqueness always means a unique
\emph{representative}, not a unique Fourier label. This establishes
genuine physical-parameter openness rather than openness only within
an equilibrium-coordinate family. No full-system invertibility or
hyperbolicity is required to continue the scalar synchronous root.

This verifies the all-pair family, universal forward witness, analytic
reverse family, rational $n=7$ witness, lower-branch family, rational
reverse-selection example, and the $n=8$ two-unstable example. The
open neighborhoods may depend on $n$. A double DM root elsewhere in
the full common domain is not claimed to support such an IFT argument.
Selection alone does not identify the first mode to cross along a
parameter path, and no path, Hopf theorem, transversality, Lyapunov
coefficient or stable periodic orbit is inferred.

The five existence categories have the exact scopes claimed:
both stable with equal or unequal representatives for every fixed size;
DM stable/MB unstable for every size; the reverse only in its proved
branch/size domains; and both unstable with distinct representatives
at the exhibited $n=8$. There is no asserted complete partition of
all physical parameter space or universal claim for every category.

\section{Independent computational corroboration and limitations}
The checker \path{A08_support/independent_check.py} was written from
the original equations and this audit's independent derivations. It
imports no P08 validation code or results. Before running, the distinct
dimensions and expected construction counts were recorded in
\path{expected_counts.json}. With
\[
 \mathcal N=\{5,6,\ldots,40,41,42,43,44,47,51\},
\]
there are 42 dimensions. The independently derived planned counts are
\[
 \sum_{n\in\mathcal N}(J(n)+1)^2=2665,\qquad
 \sum_{n\in\mathcal N}(J(n)+1)=307,\qquad
 \sum_{n\in\mathcal N}J(n)=265.
\]
These count all ordered mode pairs, strict MB cells, and positive
midpoint ties respectively. There are 12 dimensions in the reverse
family and 32 in the lower family; the union of distinct collision
offsets $\{1,n-1\}\cup\{n/2:n\text{ even}\}$ has 104 dimension-offset
pairs. These planned counts belong to the new independent checker;
reproduction of the original validation is recorded separately below.

\begin{longtable}{@{}L{0.67\textwidth}r@{}}
\toprule Independent check & Completed count\\\midrule\endhead
Universal rational witness &42\\
Every ordered stable mode pair &2665\\
Unique-branch reverse-family sizes &12\\
Lower-branch reverse-family sizes &32\\
Strict MB selection-cell constructions &307\\
Positive-midpoint MB ties &265\\
Common-domain rational witnesses, exact Fraction residuals &9\\
Common-domain equilibrium cases including repeated explicit checks &2760\\
Original-field matrix/Fourier checks, two models per collision pair &208\\
Maximum inequalities, unique/tied/all-tied maxima &312\\
Independent random physical parameter tuples &84000\\
Admissible positive-branch cases from those random tuples &40436\\
MB-stable unique-domain random cases &15741\\
\bottomrule
\end{longtable}

All 25,999 implemented assertions pass; the explicit source-label
comparison separately records E01 as a false claim, rather than
treating that discrepancy as a passing source assertion. The two
structural growth fixtures are also independently substituted in the
original fields. Exact rational arithmetic checks nine common-domain
explicit witnesses, cubic isolation signs and rational spectral bounds.
The 80-digit Decimal calculations use an independently coded cosine
Taylor series and a stated high-precision value of $\pi$, with scaled
tolerance $10^{-65}$. They are high-precision corroboration, not interval
certification. Largest scaled original equilibrium residual is below
$6.72\cdot10^{-80}$.

For collision checks a separate implementation of each original vector
field is differentiated by complex steps, preserving repeated neighbor
contributions automatically. The resulting matrices are checked by
direct action on Fourier vectors, normality and numerical eigenvalues.
The maximum scaled Fourier residual is below $4.72\cdot10^{-15}$;
the maximum scaled eigenvalue-set discrepancy is below
$7.43\cdot10^{-15}$. Full label multiplicities are supplied by the
analytic unitary diagonalization, not inferred from numerical matching.

The deterministic counterexample search uses seed 20261003 and 2,000
positive tuples per dimension, with base-10 logarithms independently
uniform in $[-3,3]$. It examines both positive DM roots when present
and MB's own positive root when $M>0$, testing the parity implications
with stated floating-point margins. It finds no parity counterexample.
It finds no random MB-stable lower-branch example; this sampling
limitation is explicit and does not weaken the exact lower-branch
construction, which is separately tested in all 32 applicable sizes.
No absence of random counterexamples is used as a proof.

The audit's exact arguments, not these finite ranges, prove all-$n$
claims, the impossibility results, openness, the blow-up formula,
normality, and $\lambda_0=-a$.

\subsection{Separate reproduction of the located historical companion}
The located companion has SHA-256
\sha{81352ce048f5bd569a4d3b803f2cbd246}{d412aa829de43f6b6d739f2521b3958}
Its archive record identifies the exact supplied P08 ZIP and all 98
final member hashes match. The inspected checker only reads its inputs
and writes its result files into the explicitly supplied output folder;
it was run into a separate temporary destination.

Before reading the saved summary or comparing a rerun, the expected
counts were independently derived from its loop structure. For the
equilibrium/matrix cases: two families across 28 sizes give 56, the
reverse and lower families add 7 and 21, the extra rational checks add
15, all ordered pairs for sizes 5--16 add 158, and distinct collision
dimension-offset cases add 14, totaling 271. Exact families count
$5+3+2+1=11$; strict cells plus ties count
$\sum_{n=5}^{32}(2J(n)+1)=280$; two models per collision pair give
28; the maximum checks give $9(4+5+6+7)=198$.

The rerun matches those independently expected counts. Its complete
summary and all per-case JSON results equal the archived originals
exactly, including the numerical maxima: matrix error
$1.888608823698555\cdot10^{-14}$, direct eigenvalue error
$2.1152019853906705\cdot10^{-14}$, and original-equilibrium residual
$1.73735320012698\cdot10^{-16}$. The inspected source, original and
reproduced summaries, reproduced results, hashes, expected counts and
comparison record are preserved under
\path{A08_support/prior_validation/}. This reproduces the historical
computation but does not turn it into mathematical proof or certify
statements it never checks. In particular its universal-witness cases
test opposite stability and a generic selection rule, without asserting
the mistaken explicit $n=8$ representative $1$.

\section{Source-range audit ledger}
The following ledger covers every line of the exact 846-line P08 file,
including headings, comments, blank lines, and the internal checklist.
PASS concerns the stated scope; MINOR ISSUE identifies required E01.
There is no MAJOR ISSUE or UNRESOLVED range. The literal example
at E01 is false, but its local correction warrants MINOR ISSUE rather
than a FALSE verdict for the valid theorem group.

\begin{longtable}{@{}L{0.10\textwidth}L{0.17\textwidth}L{0.67\textwidth}@{}}
\toprule Lines & Status & Independent disposition\\\midrule\endhead
1--13 & PASS & Current DRAFT metadata, canonical models, verified prerequisites and explicit exclusions.\\
14--33 & PASS & Same positive tuple and own equilibria; strict spectral vocabulary, full candidate labels and deferred scope.\\
34--45 & PASS & Polynomial local flow, coordinate faces and positivity only on the maximal interval.\\
46--62 & PASS & Scalar equilibria, forward growth and blow-up with correct signs and initial-data scope.\\
63--78 & PASS & Locally Lipschitz field and integrating-factor identity prove sign preservation.\\
79--93 & PASS & Synchronous lift and reciprocal solution; zero data retained and maximal intervals agree.\\
94--110 & PASS & Maximum comparison and precise global, attraction and absorption claims.\\
111--122 & PASS & Finite maximum handles ties; nonnegative inhibitory terms may be discarded; logistic comparison.\\
123--129 & PASS & Necessary conditions deduced only from synchronous examples; converses not asserted.\\
130--150 & PASS & General derivative and pre-equilibrium diagonal; coincident entries add; equilibrium used only afterward.\\
151--171 & PASS & All phases, real and imaginary signs, full spectrum, normality and negative uniform eigenvalue.\\
172--184 & PASS & Shift convention, unitary DFT and commuting-adjoint proofs are valid.\\
185--191 & PASS & Independent matrix and scalar derivations of the uniform eigenvalue.\\
192--211 & PASS & All collision formulas, Nyquist value, conjugacy and multiplicity qualifications.\\
212--228 & PASS & Strict and equality stability cases; hyperbolicity distinguished; additional sufficient global bound.\\
229--246 & PASS & Symmetric-part quadratic estimate, sum logistic bound and automatic positive MB root.\\
247--255 & PASS & Exact dissipative example with k<c; own roots and residuals; non-equivalence of sufficient bounds.\\
256--269 & PASS & Step-two quadratic, completion of square, strict concavity and constrained versus unconstrained vertex.\\
270--280 & PASS & Correct generic grid definitions; MB hypotheses are established before using P05.\\
281--297 & PASS & All strict cells, closed endpoints, positive-midpoint ties and reachable representatives.\\
298--311 & PASS & Label counts and exact selected-mode stability including uniform maximal ties.\\
312--326 & PASS & Independent gap sign and MB-specific positive-parameter realizability.\\
327--343 & PASS & Four parity classes, unattainable zero-vertex tie and amplitude-dependent gap qualification.\\
344--366 & PASS & Complete common domain, unique upper root on C-plus, and explicit branch-labelled use elsewhere.\\
367--382 & PASS & Correct imported DM coefficients and same-tuple selection controls in the q=2 comparison.\\
383--404 & PASS & Curvature, coordinate faces, uniform derivatives and whole-class dissipativity contrast with domains qualified.\\
405--414 & PASS & Three positive curvature witnesses and original residuals; no inference of MB boundedness from existence.\\
415--427 & PASS & Matching limits and fixed-source limit; derivative-sign statement read in the established q=2 scope.\\
428--437 & PASS & Every ordered pair quantified at each fixed size, with strict stable physical open sets.\\
438--455 & PASS & Strict control choices, positive-state construction and the large-k bound.\\
456--466 & PASS & Positivity, common unique-branch domain, exact original residuals and independent selection ratios.\\
467--478 & PASS & Both all-mode stability bounds; adjusting a with k preserves the chosen equilibrium and controls.\\
479--489 & PASS & Rational reverse-selection example, distinct amplitudes, modes 2/0 and both exact negative maxima.\\
490--503 & PASS & Universal same-tuple theorem; fixed-size open regions and accidental amplitude equality stated correctly.\\
504--515 & PASS & Exact original residuals, coefficients and both dispersions at the universal witness.\\
516--524 & PASS & Actual finite label J and all-size positive lower bound 3/16; nonuniformity follows from -a.\\
525--533 & PASS & Algebraic-integer exclusion of all representative ties and valid fixed-size openness.\\
534--535 & MINOR ISSUE & E01: n=8 selects MB representative 2 (labels 2 and 6), not 1; n=5,6,7 select 1.\\
536--538 & PASS & Finite Fourier label, not a continuous maximum alone, establishes instability.\\
539--549 & PASS & Normalized modal comparison and strictly positive last numerator.\\
550--557 & PASS & Exact unique-domain reverse-stability iff theorem and strict implication in other parity classes.\\
558--572 & PASS & Uniform branch stability, forced positive w and strict modal gap exclude even zero DM maxima.\\
573--589 & PASS & All-size reverse-family positivity, common-domain inequalities and original residuals.\\
590--594 & PASS & MB terminal cell and exact negative normalized maximum.\\
595--606 & PASS & DM positive terminal mode with correct 37/256 subtraction bound and 91/256 margin.\\
607--616 & PASS & DM concavity, strict terminal selection and continuation on the positive simple branch.\\
617--627 & PASS & Rational seven-node tuple, own roots, strict positivity and original residuals.\\
628--638 & PASS & Seventh-root cubic and all rational isolating signs independently verified.\\
639--654 & PASS & Every MB representative negative; exact DM positive margin; strict selection of representative 2.\\
655--665 & PASS & Separate full branch-labelled theorem with explicit lower branch and divisibility scope.\\
666--674 & PASS & Full-domain 4-divides-n obstruction and strict cosine-grid rho bound in other sizes.\\
675--693 & PASS & Lower-family coefficients, both exact positive roots, common domain and distinct branches.\\
694--710 & PASS & Positive lower uniform value, convex selection, strict all-mode MB bound and rational six-node case.\\
711--718 & PASS & Five existence categories reported without claiming each in every dimension.\\
719--738 & PASS & n=8 two-unstable example: exact residuals, distinct modes 2/1, strict vertices and positive maxima.\\
739--745 & PASS & Physical five-parameter openness at fixed size and named branch is stated precisely.\\
746--758 & PASS & Correct upper-branch scalar IFT derivative, smooth MB root and finitely many strict gaps.\\
759--765 & PASS & Correct lower-branch derivative and strict branch conditions; no uniform-in-n radius claimed.\\
766--779 & PASS & Summary preserves all theorem quantifiers and permits open MB-stable regions.\\
780--788 & PASS & Exact M=0 exponential and M=-1 blow-up witnesses; own DM roots; no MB positive root claimed.\\
789--799 & PASS & Six-term, Hopf, transversality, criticality and later proof gates explicitly deferred.\\
800--806 & PASS & Analytic checks and original-field residual identities are independently corroborated.\\
807--811 & PASS & Located companion authenticated; loop counts independently derived and the original checker/results reproduced exactly.\\
812--814 & PASS & Floating-point checks, interval certification, proof and independent audit are correctly distinguished.\\
815--846 & PASS & Checklist and scope reminders; no new theorem or later-module premise is imported.\\
\bottomrule
\end{longtable}

\section{Closure disposition and deliverables}
\textbf{Required correction: E01 only.} Correct the $n=8$ universal
witness's MB representative from $1$ to $2$, retaining the correct
$n=5,6,7$ statement. No mathematical theorem, parameter tuple,
equilibrium, spectral formula, inequality, proof construction or
branch condition needs correction. The located original validation
companion is separately reproduced and creates no additional closure
requirement. No other required change is identified.

The strongest single universal opposite-stability result independently
verified here is that the same rational tuple
$(1/4,1/8,1,1,1/8)$ has $x^*=y^*=1$, unique uniform DM selection
and $\sa(J_{\DM})=-1/4$, but unique nonuniform MB selection and
$\sa(J_{\MB})\ge3/16$ for every $n\ge5,q=2$. It persists on a
nonempty physical open set for each fixed size. Equally substantial
verified results are every ordered reachable pair with both models
stable and the exact branch-qualified reverse-stability obstructions.

\textbf{P07 clearance:} P07 is cleared after the original development
chat applies E01 to P08 and records audit closure for the corrected exact
version. It is not begun or
administratively cleared in this audit-only task. P08 remains DRAFT in
all supplied project files. No later proof or full-project audited
archive is created.

The audit-only archive is exactly
\path{cyclic_network_project_A08_audit_only.zip}. Its sole project root
is \path{cyclic_network_project/}; it includes this standalone source,
its compiled PDF, the input/member hash inventory, the independent
checker and outputs, the source-range ledger, and concise provenance
and delivery checks. It excludes P08 source, all replacement project
files, the original input archive, and build intermediates. The input
ZIP retains its original name and bytes. Compilation and visual checks
certify the reading copy, not mathematical truth.
\end{document}
